\section{Configuration Selection and Recorded Cost}
Table~\ref{tab:configuration-guide} links measured operating points to distinct objectives. These are choices within the tested settings: target and directional CLIP describe text compatibility, while retention metrics describe source similarity. CORR25--ENERGY estimates the allocation-policy contrast; its CLIP intervals span zero and do not establish semantic noninferiority.
\begin{table}[t]
\centering\small
\caption{Configuration guide for the tested operating points.}
\label{tab:configuration-guide}
\begin{tabular}{>{\raggedright\arraybackslash}p{0.27\columnwidth}>{\raggedright\arraybackslash}p{0.63\columnwidth}}
\toprule
Objective & Measured choice and evidence\\
\midrule
Retention within correction arms & U25: best mean retention metrics in this family, 91 calls.\\
Target-text compatibility & BOTH25: highest mean target CLIP, 0.2733; 97 calls.\\
Directional response and call budget & FlowEdit N26: directional CLIP 0.1177 and 52 calls.\\
Allocation attribution & CORR25 versus ENERGY: small LPIPS gain under common-state calibration, with measured trajectory drift.\\
\bottomrule
\end{tabular}
\end{table}

RF-Edit I08 is a high-retention, weak-editing working point: its factual-edit count is 4/40 and edit-and-preserve count 3/40. Its retention maxima are shown in the complete tables alongside that editing response.

Table~\ref{tab:runtime-cost} summarizes recorded synchronized generation time and CUDA allocation/reservation peaks. These are arithmetic means over native outputs (80 per arm, 40 for RF-Edit), with evaluation time excluded. They describe this sequential run, including its cache and startup state; the arms were not randomized for a dedicated timing benchmark. CUDA peaks measure allocator usage, not total device memory. NFE remains the standalone structural call count.
\begin{table}[t]
\centering\small
\caption{Recorded mean generation time and allocator peaks. GiB uses $2^{30}$ bytes.}
\label{tab:runtime-cost}
\begin{tabular}{lrrr}
\toprule
Arm & Seconds & Alloc. GiB & Reserv. GiB\\
\midrule
CORR25 & 66.01 & 22.40 & 22.52\\
BOTH25 & 55.10 & 22.40 & 22.52\\
ENERGY & 54.94 & 22.40 & 22.52\\
U25\_KV & 53.23 & 22.40 & 22.52\\
FlowEdit N26 & 39.64 & 22.33 & 22.51\\
RF-Edit I08 & 54.97 & 22.33 & 22.50\\
ETA03 & 44.65 & 22.34 & 22.51\\
\bottomrule
\end{tabular}
\end{table}

